\lab{8}{Булев поиск}
\task
Реализовать запросы с И (пробел или \code{&&}), ИЛИ (\code{||}),
НЕ (\code{!}) и скобками; CLI с запросами по строкам stdin и веб-форму
с заголовками, URL и постраничной выдачей. Проверить корректность
и обсудить сложные запросы и скорость.

\method
Рекурсивный парсер строит дерево запроса. Приоритеты:
\code{!}, затем \code{&&}, затем \code{||}. Скобки задают другой
порядок, соседние атомы в строгом режиме соединяются операцией И.
Пустой запрос даёт пустое множество. Одиночные \code{&}, незакрытые
скобки и незавершённая операция считаются синтаксической ошибкой.

Терм нормализуется так же, как при индексации. Его постинг-лист читается
из файла; результат поддерева - отсортированное множество \(docID\).
Пересечение и объединение выполняются двумя указателями за
\(O(a+b)\), отрицание - дополнением относительно диапазона всех
документов \([0,N)\). Поэтому пустой документ может входить в ответ
на \code{!memory}. Повторы идентификаторов в выдаче исключаются.

CLI загружает индекс и принимает несколько запросов за запуск.
Ответ JSON содержит общее число совпадений, смещение, время,
число прочитанных постингов и результаты. Веб-обвязка запускает
локальный поиск и экранирует заголовки и ссылки при формировании HTML.
Параметры \code{search_limit=50} и \code{search_offset} реализуют
постраничное извлечение. Веб-сервер запускается отдельно от индексатора.

\journal
Разбор запроса отделён от операций над множествами. Это позволило
позднее сохранить язык булевых операций при переходе к координатному
индексу. Добавлены обработка переменного количества пробелов,
двойное отрицание и диагностика некорректного синтаксиса.
Нормализация запросов перенесена в общий конвейер токенизации,
чтобы запрос не расходился с представлением документа.

\testing
На искусственном корпусе из четырёх документов проверяются все
операции, приоритеты, пустой документ, неизвестный терм и пустой ввод.
Для \code{memory || connections} ожидается объединение без повторов;
для \code{connections !mutex} - исключение документа с mutex.
Дерево \code{alpha || beta && !delta} должно иметь корневую операцию
ИЛИ и правую ветвь И. Два модульных теста пройдены.

Для веб-интерфейса дополнительно требуются ручные проверки пустой
выдачи, перехода к следующей странице, спецсимволов в заголовке
и ошибки запроса. Эти проверки не заменяются C++-тестами парсера.

\results
Примеры запросов для демонстрации:
\begin{lstlisting}
memory allocation
memory && (allocation || mapping)
(mutex || spinlock) && !futex
!memory
\end{lstlisting}
\begin{lstlisting}
cd 8_boolean
just build
just search < queries.txt
just web
\end{lstlisting}
Готовый индекс уже существует; повторное \code{build-index} для поиска
не требуется. Время работы зависит от количества прочитанных
постингов и размеров промежуточных множеств. Замеры координатного
движка из следующей работы не характеризуют скорость этой реализации.

Дорогие запросы - объединения частых слов, чистое отрицание и широкие
подвыражения, создающие промежуточные списки почти из всех документов.
Выдача только первых 50 заголовков не устраняет стоимость полного
вычисления множества совпадений в текущем алгоритме.

\conclusions
Язык запросов позволяет точно задавать условия присутствия слов,
но порядок выдачи по \(docID\) слабо связан с полезностью документа.
Возможные улучшения - выполнение И начиная с самого короткого
списка, курсоры вместо материализации промежуточных множеств,
прыжки и ранжирование. Эти направления реализуются в следующих работах.
